% 分类目录：dl
\documentclass[12pt, a4paper]{article}
\usepackage[margin=2.5cm]{geometry}
\usepackage{ctex}
\usepackage{amsmath, amssymb}
\usepackage{listings}
\usepackage{xcolor}
\usepackage{tabularx}
\usepackage{hyperref}

\lstset{
    language=Python,
    basicstyle=\ttfamily\small,
    keywordstyle=\color{blue},
    commentstyle=\color{green!50!black},
    stringstyle=\color{red},
    showstringspaces=false,
    breaklines=true,
    numbers=left,
    numberstyle=\tiny\color{gray},
    frame=single
}

\title{知识点精讲：深度学习模型优化算法 (Optimization Algorithms)}
\author{}
\date{}

\begin{document}
\maketitle

\section{优化问题的数学本质}
在深度学习中，优化目标是寻找参数 $\theta$ 使得损失函数 $J(\theta)$ 最小。在高维非凸空间中，优化面临的核心挑战如下：

\begin{itemize}
    \item \textbf{鞍点 (Saddle Points)}：梯度为 0 但非极值的点。在高维空间中，鞍点数量远超局部极小值。判断准则是 Hessian 矩阵的特征值同时包含正数和负数。
    \item \textbf{病态曲率 (Ill-conditioned Curvature)}：Hessian 矩阵的条件数 (Condition Number) 过大，导致梯度方向并不指向局部极小值点，使得优化过程在陡峭方向震荡，在平滑方向进展缓慢。
\end{itemize}

\section{经典梯度下降及其变体}

\subsection{动量法 (Momentum)}
引入一阶矩（动量 $v$）来加速收敛并抑制震荡：
$$ v_t = \beta v_{t-1} + (1-\beta) \nabla J(\theta_t) $$
$$ \theta_t = \theta_{t-1} - \eta v_t $$
\textbf{直觉}：模拟物理惯性，在梯度不一致的方向相互抵消，在一致的方向叠加。

\subsection{Nesterov 加速梯度 (NAG)}
在计算梯度前，先根据当前动量向前“展望”一步：
$$ g_t = \nabla J(\theta_{t-1} - \eta \beta v_{t-1}) $$
$$ v_t = \beta v_{t-1} + (1-\beta) g_t, \quad \theta_t = \theta_{t-1} - \eta v_t $$
\textbf{核心}：NAG 具有“提前刹车”的能力，在坡度改变前预感知，从而提高稳定性。

\section{自适应学习率算法}

\subsection{Adam (Adaptive Moment Estimation)}
结合了 Momentum（一阶矩 $m$）和 RMSProp（二阶矩 $v$）：
$$ m_t = \beta_1 m_{t-1} + (1-\beta_1)g_t $$
$$ v_t = \beta_2 v_{t-1} + (1-\beta_2)g_t^2 $$
\textbf{偏置修正 (Bias Correction)}：由于 $m_0, v_0$ 初始化为 0，初期估计会严重向 0 偏移：
$$ \hat{m}_t = \frac{m_t}{1-\beta_1^t}, \quad \hat{v}_t = \frac{v_t}{1-\beta_2^t} $$
$$ \theta_t = \theta_{t-1} - \eta \frac{\hat{m}_t}{\sqrt{\hat{v}_t} + \epsilon} $$

\subsection{AdamW：权重衰减的解耦}
在 Adam 中，L2 正则化项合并在梯度中会受到二阶矩自适应缩放的影响，导致正则化效果变差。AdamW 将权重衰减直接应用在参数更新步，修正了这一逻辑。

\section{现代 SOTA 优化器}

\subsection{Lion (EvoLved Sign mOmeNtum)}
Google 通过进化算法发现，仅使用梯度的符号位即可获得极佳性能：
$$ u_t = \text{sign}(\beta_1 m_{t-1} + (1-\beta_1)g_t) $$
$$ \theta_t = \theta_{t-1} - \eta u_t, \quad m_t = \beta_2 m_{t-1} + (1-\beta_2)g_t $$
\textbf{优势}：显存占用比 AdamW 少 50\%（不存储二阶增量），且在 Transformer 任务中具有更强的泛化性。

\subsection{Sophia：轻量级二阶优化器}
针对大语言模型 (LLM) 设计。利用对 Hessian 矩阵对角元的随机估计进行剪切优化：
$$ \theta_t = \theta_t - \eta \cdot \text{clip}\left(\frac{m_t}{\max(h_t, \gamma)}, \rho\right) $$
在同等训练开销下，比 AdamW 同步快 2 倍以上。

\section{泛化策略：平坦极小值 (Flat Minima)}
\textbf{Lookahead}：维护一组慢权重，定期同步快权重的平均方向，步步为营。
\textbf{SWA (Stochastic Weight Averaging)}：在训练后期对多个模型快照求算术平均，旨在寻找损失曲面上更加“宽阔平坦”的区域。

\section{核心代码参考 (PyTorch)}
\begin{lstlisting}
import torch.optim as optim

# 1. 工业级 AdamW 配置
optimizer = optim.AdamW(model.parameters(), lr=1e-3, weight_decay=0.01)

# 2. 结合 CosineAnnealing 调度器
scheduler = optim.lr_scheduler.CosineAnnealingLR(optimizer, T_max=100)

# 3. 学习率 Warmup (Linear Warmup)
# 通常在大规模训练初期手动实现或使用工具库实现
\end{lstlisting}

\end{document}
